\documentclass[11pt]{article}
\usepackage[margin=0.9in]{geometry}
\usepackage{amsmath,amssymb,amsthm,mathtools}
\usepackage[T1]{fontenc}
\usepackage{lmodern}
\usepackage[hidelinks]{hyperref}
\newtheorem{theorem}{Theorem}
\newtheorem{lemma}[theorem]{Lemma}
\newtheorem{proposition}[theorem]{Proposition}
\newtheorem{corollary}[theorem]{Corollary}
\theoremstyle{remark}
\newtheorem{remark}[theorem]{Remark}
\DeclareMathOperator{\Mat}{M}
\DeclareMathOperator{\Pic}{Pic}
\DeclareMathOperator{\im}{im}
\DeclareMathOperator{\Spec}{Spec}
\title{Primitive idempotents that cannot be homogenized in Koszul algebras}
\author{Henry Zweiman}
\date{September 9, 2026}
\begin{document}
\maketitle
\begin{abstract}
We construct primitive idempotents that no algebra automorphism can carry to a homogeneous idempotent in a locally finite standard graded algebra with semisimple degree zero. The algebra is a matrix algebra over a two-dimensional domain defined by two quadrics, and is Koszul. This gives a negative answer to Dramburg's Question 6.2, including over algebraically closed fields of characteristic zero. The same fixed algebra contains such an idempotent congruent to a constant matrix modulo any prescribed power of the positive-degree ideal. Each is conjugate to that constant matrix after completion. The proof combines an explicit matrix with a line-module obstruction, detected by its two trivializations over a Veronese ring and a polynomial quotient. Counterexamples exist over every field, in matrix size two outside characteristic two and size three in characteristic two.
\end{abstract}

\section{The question and the construction}

A nonnegatively graded $k$-algebra $A=\bigoplus_{j\geq0}A_j$ is standard graded if it is generated by $A_0$ and $A_1$, locally finite if each $A_j$ is finite dimensional, and semiconnected if $A_0$ is semisimple. Dramburg asks whether every complete set of primitive orthogonal idempotents of such an algebra can be carried to its degree-zero parts by an automorphism \cite[Question 6.2]{Dramburg}. His earlier question explicitly includes the possibility of assuming that $k$ is algebraically closed of characteristic zero \cite{MO}. We give counterexamples satisfying all these conditions, with the additional Koszul property.

The relation between idempotent matrices and projective modules is classical. Petersson describes line-module obstructions to automorphisms of two-by-two matrix algebras \cite[Sections 5--6]{Petersson}. Kanwar, Leroy and Matczuk study conjugacy in polynomial and power-series extensions, and give a quaternion-polynomial example using Ojanguren--Sridharan \cite[p.~9 of the author preprint]{KLM}. De Jong also describes nonhomogeneous idempotents arising from line bundles in a Laurent extension \cite{deJong}. Our contribution is an explicit split construction over every ground field, with a quadratic complete-intersection center, and an arbitrarily deep finite-jet failure in one fixed Koszul algebra. The underlying projective-module method is not new.

Fix a field $k$ and put
\begin{equation}\label{eq:ring}
 R=k[a,b,c,z]/(b^2-ac,\ c^2-az),\qquad
 \deg a=\deg b=\deg c=\deg z=1.
\end{equation}
For $r\geq0$, define four elements of $R$ by
\begin{equation}\label{eq:powers}
 u_{2,r}=cz^{2r+1},\quad u_{3,r}=bz^{3r+2},\quad
 u_{4,r}=az^{4r+3},\quad u_{5,r}=bcz^{5r+3},
\end{equation}
and set
\begin{equation}\label{eq:matrix}
 p_r=\begin{pmatrix}
 1-u_{4,r}&u_{2,r}+u_{3,r}+u_{4,r}+u_{5,r}\\
 u_{2,r}-u_{3,r}&u_{4,r}
 \end{pmatrix}.
\end{equation}
Write $E_{ij}$ for the standard matrix units. In $A=\Mat_n(R)$ let
\[
 e_r=\operatorname{diag}(p_r,0_{n-2}),\qquad
 f_r=\operatorname{diag}(1-p_r,0_{n-2}),\qquad
 A_j=\Mat_n(R_j),\quad A_+=\bigoplus_{j>0}A_j.
\]
The notation for the zero block is omitted when $n=2$.

\begin{theorem}\label{thm:main}
Let $n\geq2$ be an integer whose image in $k$ is nonzero. For the ring and matrices above, the following statements hold.
\begin{enumerate}
\item $R$ is a two-dimensional standard graded domain and a quadratic complete intersection. The algebra $A$ is locally finite, standard graded and Koszul, with $A_0=\Mat_n(k)$.
\item For every $r\geq0$, the elements $e_r,f_r,E_{33},\ldots,E_{nn}$ form a complete set of primitive orthogonal idempotents. Their degree-zero parts are $E_{11},E_{22},E_{33},\ldots,E_{nn}$.
\item No $k$-algebra automorphism of $A$ carries $e_r$ to any homogeneous idempotent. In particular, no automorphism carries the displayed complete system to its degree-zero parts, even allowing a permutation.
\item We have $e_r-E_{11}\in A_+^{2r+2}$. Nevertheless $e_r$ and $E_{11}$ are conjugate by a unit congruent to $1$ modulo $A_+^{2r+2}$ in the $A_+$-adic completion. Their images are conjugate in every quotient $A/A_+^N$, $N\geq1$.
\end{enumerate}
\end{theorem}

Taking $n=2$ when $\operatorname{char}k\ne2$, and $n=3$ otherwise, gives a counterexample over every field. In particular, $n=2$ works over $\mathbb C$. The theorem concerns arbitrary algebra automorphisms, not just conjugation by units or automorphisms fixing the center pointwise.

\section{The center and the auxiliary rings}

\begin{lemma}\label{lem:ring}
The map in \eqref{eq:ring} induced by
\[
 a\longmapsto s^4,\qquad b\longmapsto s^3t,\qquad
 c\longmapsto s^2t^2,\qquad z\longmapsto t^4
\]
identifies $R$ with the indicated subring of $k[s,t]$. The two defining quadrics form a Gr\"obner basis. The ring $R$ is Koszul.
\end{lemma}
\begin{proof}
Use lexicographic order $b>c>a>z$. The leading monomials $b^2$ and $c^2$ are relatively prime, so Buchberger's criterion gives the asserted Gr\"obner basis. Its standard monomials are
\begin{equation}\label{eq:standard}
 a^i z^j b^\epsilon c^\delta,
 \qquad i,j\geq0,\quad \epsilon,\delta\in\{0,1\}.
\end{equation}
Their images in $k[s,t]$ have exponent pairs
\[
 (4i+3\epsilon+2\delta,\ 4j+\epsilon+2\delta).
\]
The first exponent modulo $4$ determines $(\epsilon,\delta)$, since the four possibilities are $0,3,2,1$. The two exponents then determine $i,j$. Thus the standard monomials have distinct images, proving injectivity and the domain assertion.

Formula \eqref{eq:standard} also makes $R$ a free module of rank four over $k[a,z]$, with basis $1,b,c,bc$. Consequently $\dim R=2$, and the defining ideal has height two. Two generators of height two in a polynomial ring form a regular sequence, so this is a complete intersection. Finally, a standard graded algebra with a quadratic Gr\"obner basis is Koszul; see, for example, \cite[Section 1]{Conca}.
\end{proof}

Work henceforth in $k[s,t]$, using the identifications of Lemma~\ref{lem:ring}. Introduce
\begin{equation}\label{eq:auxiliary}
 B=k[s^4,s^3t,s^2t^2,st^3,t^4],\qquad d=st^3,\qquad C=k[z].
\end{equation}
Thus $B$ is the fourth Veronese subring. Give its degree-four monomials degree one. Let $J$ be the $k$-span in $B$ of all its monomials having $s$-exponent at least two. This is an ideal. In the quotient only the monomials $z^j$ and $dz^j$, $j\geq0$, remain. Hence
\begin{equation}\label{eq:diagram}
 D=B/J\cong k[z,d]/(d^2).
\end{equation}
The symbol $d$ in $D$ denotes the residue class of $st^3$.

The inclusion $R\to B$ followed by the quotient to $D$ sends $a,b,c$ to zero and $z$ to $z$. It therefore agrees with the composition $R\to C\to D$, where $R\to C$ kills $a,b,c$. We will use this commutative square of ring maps. We do not need flatness of $B$ over $R$, a computation of $\Pic(B)$, or a description of the full normalization of $R$.

Both $B$ and $C$ have unit group $k^\times$. For $B$, this follows because it is a connected positively graded domain: the highest nonzero degrees of two nonzero elements add under multiplication, so a product equal to $1$ forces both elements to have degree zero. The assertion for $C=k[z]$ is immediate.

\section{An explicit line-module obstruction}

For an indeterminate $x$, consider the matrix
\begin{equation}\label{eq:T}
 T(x)=\begin{pmatrix}
 1+x+x^2+x^3&-x^2\\
 x^2&1-x
 \end{pmatrix}\in\Mat_2(\mathbb Z[x]).
\end{equation}
Its determinant is $(1-x)(1+x+x^2+x^3)+x^4=1$. Direct multiplication gives
\begin{equation}\label{eq:conjugate}
 T(x)E_{11}T(x)^{-1}
 =\begin{pmatrix}1-x^4&x^2+x^3+x^4+x^5\\x^2-x^3&x^4\end{pmatrix}.
\end{equation}
Put $\xi_r=dz^r\in B$. The identities
\[
 \xi_r^2=cz^{2r+1},\quad \xi_r^3=bz^{3r+2},\quad
 \xi_r^4=az^{4r+3},\quad \xi_r^5=bcz^{5r+3}
\]
show that \eqref{eq:conjugate}, evaluated at $\xi_r$, is precisely $p_r$ and has entries in $R$. Since $R\to B$ is injective, $p_r$ is idempotent over $R$.

Set $L_r=p_rR^2$. The determinant and trace of $p_r$ are $0$ and $1$. Over every residue field its image has dimension one; thus $L_r$ and $(1-p_r)R^2$ are finitely generated projective modules of rank one. Here a line module means a projective module of rank one.

\begin{proposition}\label{prop:line}
For every $r\geq0$ and every positive integer $m$ nonzero in $k$, the line module $L_r^{\otimes m}$ is not free. The line modules $L_r$, $r\geq0$, are pairwise nonisomorphic.
\end{proposition}
\begin{proof}
After base change to $B$, the module $L_r$ has the basis vector
\[
 v_r=T(\xi_r)\binom10.
\]
Indeed, conjugation in \eqref{eq:conjugate} identifies the image with the span of the first column of an invertible matrix. After base change to $C$, the idempotent becomes $E_{11}$, so $L_r\otimes_R C$ has basis $w=\binom10$.

These identifications are compatible over $D$. Since $\xi_r^2=0$ in $D$, we have
\begin{equation}\label{eq:transition}
 T(\xi_r)\bmod J=\operatorname{diag}(1+dz^r,1-dz^r),
 \qquad v_r=(1+dz^r)w\quad\hbox{over }D.
\end{equation}
Tensor powers of these basis vectors consequently differ by $(1+dz^r)^m$.

Suppose $L_r^{\otimes m}\cong R$ and choose a global basis. Over $B$ its expression is $\beta v_r^{\otimes m}$ for some $\beta\in B^\times$; over $C$ it is $\gamma w^{\otimes m}$ for some $\gamma\in C^\times$. These scalars must be units because both expressions are bases. They lie in $k^\times$, by the unit computations above. Compatibility over $D$ gives
\[
 \beta(1+dz^r)^m=\gamma.
\]
The left side is $\beta(1+mdz^r)$. The elements $1$ and $dz^r$ are linearly independent over $k$, and $m\ne0$ in $k$, a contradiction.

Likewise, an isomorphism $L_r\cong L_s$ becomes multiplication by units in the two chosen trivializations. Compatibility implies an equality
\[
 \beta(1+dz^r)=\gamma(1+dz^s),\qquad \beta,\gamma\in k^\times.
\]
Comparing the constant terms gives $\beta=\gamma$, and then $z^r=z^s$ in $k[z]$, so $r=s$.

All base changes used here are applied to direct summands of free modules, and to a hypothesized isomorphism. They require no flatness assumption.
\end{proof}

In characteristic zero, each $[L_r]$ has infinite order in $\Pic(R)$. We have not computed this Picard group, and need no assertion about its other elements.

\section{Excluding arbitrary automorphisms}

The following matrix-unit argument is a useful form of the classical projective-module obstruction. It avoids an assumption that an automorphism fixes the center.

\begin{lemma}\label{lem:automorphism}
Let $S$ be a commutative ring and let $e\in\Mat_n(S)$ be an idempotent whose image $P=eS^n$ is a line module. If a ring automorphism of $\Mat_n(S)$ carries $e$ to $E_{11}$, then $P^{\otimes n}\cong S$.
\end{lemma}
\begin{proof}
Let $\alpha$ be such an automorphism and set $q_{ij}=\alpha^{-1}(E_{ij})$. These matrices satisfy the matrix-unit identities, and $q_{11}=e$. They yield a direct sum
\[
 S^n=\bigoplus_{i=1}^n q_{ii}S^n.
\]
Multiplication by $q_{i1}$ is an $S$-linear isomorphism from $P$ to $q_{ii}S^n$, with inverse multiplication by $q_{1i}$. Thus $S^n\cong P^{\oplus n}$. Taking the top exterior power gives $S\cong P^{\otimes n}$.

The matrices $q_{ij}$ act $S$-linearly regardless of the action of $\alpha$ on scalar matrices. In particular, the proof applies to automorphisms acting nontrivially on the center.
\end{proof}

\begin{proof}[Proof of Theorem~\ref{thm:main}]
The ring statements follow from Lemma~\ref{lem:ring}. The entrywise grading on $A$ is standard and locally finite, and $A_0=\Mat_n(k)$ is semisimple. For the Koszul assertion, take a linear graded projective resolution of $k$ over $R$ and apply the exact graded Morita functor
\[
 \Mat_n(R)E_{11}\otimes_R -.
\]
It gives a linear resolution of the simple $\Mat_n(k)$-module $k^n$, regarded as an $A$-module. Taking $n$ copies gives a linear resolution of $A_0$ as a left $A$-module, which is the Koszul condition with semisimple degree zero.

The identities for $p_r$ give the stated orthogonal decomposition of $1$. Every displayed summand has image of rank one over $R$. For a line module $P$ one has $\operatorname{End}_R(P)\cong R$, as can be checked after localization. Since $R$ is a domain, this endomorphism ring has no nontrivial idempotents. Hence each displayed matrix idempotent is primitive. The degree-zero assertions follow directly from \eqref{eq:powers}.

The image of $e_r$ is isomorphic to $L_r$. Proposition~\ref{prop:line} says that $L_r^{\otimes n}$ is not free. Lemma~\ref{lem:automorphism} therefore excludes an automorphism taking $e_r$ to $E_{11}$. A nonzero homogeneous idempotent in a nonnegative grading has degree zero, since its square has twice its degree. A primitive idempotent in $\Mat_n(k)$ has rank one: any larger rank splits into orthogonal constant idempotents in $A$. Every constant rank-one idempotent is conjugate over $k$ to $E_{11}$. Composing with that constant conjugation proves the third assertion, including the statement about permutations of the complete system.

The four elements $u_{j,r}$ have respective degrees $j(r+1)$ for $2\leq j\leq5$. Thus $e_r-E_{11}\in A_{\geq2r+2}=A_+^{2r+2}$, where equality uses the standard grading. Put $E=E_{11}$ and
\begin{equation}\label{eq:formal}
 W=e_rE+(1-e_r)(1-E)=1+(e_r-E)(2E-1).
\end{equation}
Then $WE=e_rW$, by idempotence. Also $W-1\in A_+^{2r+2}$. Therefore the geometric series for $(1+(W-1))^{-1}$ converges in the $A_+$-adic completion, and becomes a finite inverse in every quotient $A/A_+^N$. In each such ring $e_r=WEW^{-1}$, proving the final assertion.
\end{proof}

\begin{corollary}\label{cor:jets}
For the fixed algebra $A$ in Theorem~\ref{thm:main} and every $N\geq1$, there is a primitive idempotent congruent to $E_{11}$ modulo $A_+^N$ whose orbit under all algebra automorphisms contains no homogeneous idempotent. The family $e_r$ consists of pairwise distinct inner-conjugacy classes.
\end{corollary}
\begin{proof}
Choose $r$ with $2r+2\geq N$. For the last assertion, an inner conjugation between $e_r$ and $e_s$ would induce an $R$-linear isomorphism between their images $L_r$ and $L_s$, contradicting Proposition~\ref{prop:line} when $r\ne s$.
\end{proof}

\section{Scope and consequences}

The construction gives a negative answer to the homogenization question even for a prime algebra finite over its center, with center a two-dimensional quadratic complete-intersection domain. The obstruction is algebraic: completion removes it, and no fixed finite order of agreement with the degree-zero part detects it. This does not assert that the $e_r$ lie in distinct orbits under all automorphisms; we only need to exclude the homogeneous orbit for each of them.

There is no conflict with the uniqueness theorem for standard semiconnected locally finite gradings \cite[Theorem 1.2]{Dramburg}. That theorem compares gradings already given on an algebra. The primitive systems here cannot be made homogeneous by choosing another grading of that class: a graded isomorphism from such a grading to the present one would carry $e_r$ to a homogeneous idempotent, contrary to Theorem~\ref{thm:main}.

The center in \eqref{eq:ring} is not seminormal: $d$ belongs to its fraction field, since $d=bc/a$, and $d^2,d^3\in R$ while $d\notin R$. This failure provides the nonconstant nilpotent unit $1+dz^r$ in the comparison ring $D$. The example therefore does not settle analogous questions restricted to normal centers or to algebras of finite global dimension. These additional restrictions require separate arguments.

\begin{thebibliography}{9}
\bibitem{Conca}
A. Conca, \emph{Koszul algebras and Gr\"obner bases of quadrics}, arXiv:0903.2397v1 (2009), Section 1. \url{https://arxiv.org/abs/0903.2397}.

\bibitem{deJong}
A. J. de Jong, \emph{Graded idempotents}, Stacks Project Blog, September 9, 2013. \url{https://www.math.columbia.edu/~dejong/wordpress/?p=3554}.

\bibitem{Dramburg}
D. Dramburg, \emph{Isomorphisms of graded semiconnected algebras}, arXiv:2609.03288v1 (2026), especially Definition 1.1, Theorem 1.2 and Question 6.2. \url{https://arxiv.org/abs/2609.03288}.

\bibitem{MO}
D. Dramburg, \emph{Are idempotents in a nonnegatively graded algebra conjugate to homogeneous idempotents?}, MathOverflow question 479788, September 29, 2024. \url{https://mathoverflow.net/q/479788}.

\bibitem{KLM}
P. Kanwar, A. Leroy and J. Matczuk, \emph{Idempotents in ring extensions}, Journal of Algebra \textbf{389} (2013), 128--136. \url{https://doi.org/10.1016/j.jalgebra.2013.05.010}.

\bibitem{Petersson}
H. P. Petersson, \emph{Idempotent 2-by-2 matrices}, author preprint, Sections 5--6. \url{https://www.fernuni-hagen.de/mi/fakultaet/emeriti/docs/petersson/idem22.pdf}.
\end{thebibliography}
\end{document}
